\section{Sharing intermediate vertices in an order-five jet construction}

\subsection*{Finding M-015: order-five jets at matrix size \(r^2\)}
Let \(r\geq4\) and \(n\geq r^2\). The map from
\(n\)-by-\(n\) linear pencils in \(r\) variables to their
determinant coefficient jets through degree five is dominant.
It has an explicit integral point with a unit Jacobian minor.
Consequently
\begin{equation}\label{eq:M015K}
 \lambda\vdash nd,\quad\ell(\lambda)\leq r,\quad\lambda_2\leq5
       \quad\Longrightarrow\quad K_n(\lambda)=a_\lambda(d[n]).
\end{equation}
In particular an actual internally padded permanent obstruction
at \(n\geq m^4\), \(m\geq2\), requires
\begin{equation}\label{eq:M015per}
                              \lambda_2\geq6.
\end{equation}
The earlier IP bound assumed \(n>3m^4\). This result
removes that factor for the second-row conclusion alone;
the other IP staircase inequalities retain their own hypotheses.

\paragraph{The pencil and its core.}
Set \(R=r-1\geq3\), write the variables as
\(x,z_1,\ldots,z_R\), and let
\(\sigma=(1\ 2\ \cdots\ R)\). The directed graph has
the following distinct vertices:
\[
 0;\qquad v_i\ (1\leq i\leq R);\qquad
 p_{ij},q_{ij}\ (1\leq i\leq j\leq R).
\]
Its size is \(1+R+2\binom{R+1}{2}=(R+1)^2=r^2\).
Its edges, and no others, are
\[
 \begin{array}{c|c}
 \text{edge}&\text{label}\\\hline
 p_{ij}\longrightarrow v_i&z_j\\
 v_i\longrightarrow0&z_i\\
 0\longrightarrow v_i&z_{\sigma(i)}\\
 v_{\sigma^{-1}(j)}\longrightarrow q_{ij}&z_i.
 \end{array}                                             \tag{M-015.1}
\]
Let \(B\) be its labeled adjacency matrix and put
\(A=xI-B\). Unlike M-013's acyclic graph, this graph has
two-edge cycles through the central vertex. The inward leaves
\(p_{ij}\) have no incoming edges, and the outward leaves
\(q_{ij}\) have no outgoing edges. At \(x=1\), therefore,
\[
        \det A=D:=1-\sum_{i=1}^R z_i z_{\sigma(i)}.
                                                               \tag{M-015.2}
\]
Indeed only the central star contributes nontrivially to
the determinant; its Schur complement is \(D\).

Let \(C\) be the \((R+1)\)-vertex star block at \(x=1\).
The elementary block inverse gives
\[
 \begin{split}
 D(C^{-1})_{00}&=1,\\
 D(C^{-1})_{0k}&=z_{\sigma(k)},\\
 D(C^{-1})_{ik}&=D\delta_{ik}+z_i z_{\sigma(k)}.
 \end{split}                                             \tag{M-015.3}
\]
These rational identities are cofactor polynomial identities
after multiplication by \(D\), so they apply without
restricting to points where \(D\ne0\).
The two leaf edges also give
\[
 D(A^{-1})_{p_{ij},q_{ab}}
     =z_jz_a\left(D\delta_{i,\sigma^{-1}(b)}+z_i z_b\right).
                                                               \tag{M-015.4}
\]
Here \(A^{-1}\) means the inverse after \(x=1\).

\paragraph{Pure monomial cofactors.}
Equation (M-015.3) supplies the constant and every linear
monomial. It also supplies every quadratic monomial from
an off-diagonal entry. For \(z_a z_b\), use row \(a\)
and column \(\sigma^{-1}(b)\) if these differ. Otherwise
reverse \(a,b\). Both choices could fail only if
\(b=\sigma(a)\) and \(a=\sigma(b)\), which would be
a fixed point or two-cycle of \(\sigma\). Neither exists
because \(R\geq3\).

Every monomial of degree four can be written
\(z_i z_j z_a z_b\) with \(i\leq j\), \(a\leq b\)
and \(b\ne\sigma(i)\). Here is a complete check.
If its largest and smallest indices differ by at least two,
put a smallest index into the first pair and a largest into
the second. The first minimum is then \(i\), the second
maximum is \(b\geq i+2\), and \(b\ne\sigma(i)\).
If every index is the same, \(b=i\ne\sigma(i)\).
The remaining case uses just the adjacent indices \(t,t+1\).
For their multiplicities \((3,1)\), take pairs
\((t,t+1)\), \((t,t)\); for \((1,3)\), take
\((t+1,t+1)\), \((t,t+1)\); and for \((2,2)\), take
\((t+1,t+1)\), \((t,t)\). The first two choices use
absence of fixed points. The last uses
\(\sigma(t+1)\ne t\), including the wraparound case
\(t+1=R\), since then \(R-1\ne1\).
These are all degree-four multisets. Equation (M-015.4)
therefore supplies each degree-four monomial as a pure
cofactor with coefficient one.

\begin{proof}[Proof of jet dominance and the multiplicity conclusions]
Varying entry \((u,v)\) of the pencil by a linear form
\(H\) differentiates its determinant to
\(H(1,z)\det A(1,z)(A(1,z)^{-1})_{vu}\).
The pure cofactors above have degrees zero, one, two and
four. Taking \(H=x\) yields those degrees directly;
taking \(H=z_i\) also yields every monomial of degrees
three and five. Thus the Jacobian of the projection onto
all coefficients through degree five is surjective at this
integral pencil. Each chosen elementary coefficient direction
maps to one distinct monomial with coefficient one, so a
Jacobian minor is the identity over \(\mathbb Z\).
For \(n>r^2\), pad by \(xI_{n-r^2}\); after \(x=1\)
the same cofactor directions remain.

Translation by the coefficient vector \(D\) and restriction
to these parameter directions give a polynomial map
\(F(t)=t+O(t^2)\) over \(\mathbb Z\). As in M-013,
the lowest homogeneous term of any nonzero polynomial
survives composition with \(F\). The pullback is injective
after every commutative base change, and the coefficient
map is dominant over \(\mathbb C\).

The support theorem F-016 places every highest-weight
function of second row at most five in this coefficient
algebra. A determinant-ideal vector in that algebra vanishes
after the dominant projection and must be zero. This holds
for the complete highest-weight space, proving
\eqref{eq:M015K} in characteristic zero.

For \eqref{eq:M015per}, a positive permanent multiplicity
has length \(r\leq m^2\). If \(r\geq4\), then
\(n\geq m^4\geq r^2\), and the preceding conclusion
excludes \(\lambda_2\leq5\). If \(r\leq3\), the
full determinant multiplicity results O-006--O-007 already
exclude every actual obstruction under these hypotheses.
This proves the asserted necessary condition.
\end{proof}

\paragraph{Grothendieck dictionary and method.}
The matrix and all cofactor identities are defined over
\(\mathbb Z\). The target is the finite free coefficient
module of \(\mathbb Z[z]/(z)^6\), and the unit Jacobian
minor gives an actual universal injection on polynomial
coordinate algebras. The coefficient weights are the torus
comodules of SGA~3, I, Proposition~4.7.3, pp.~35--36;
the whole Schur-copy support is F-016's use of SGA~6, VI,
Theorem~3.3, pp.~392--393, and Macdonald, I(4.3), p.~63,
Appendix~A(5.3)--(5.6), pp.~155--156.
Sharing intermediate vertices introduces cycles and changes
the base determinant to \(D\). The explicit off-diagonal
cofactor calculation, rather than acyclicity or an exact
truncated-form representation, supplies dominance. Higher
coefficients remain free. The conclusion is an ideal-comodule
vanishing statement over characteristic zero, not a virtual
\(K_0\) equality or a saturation inference.
